\documentclass[UTF8,fontset=fandol]{ctexart}
\usepackage[a4paper,margin=21mm]{geometry}
\usepackage{amsmath,amssymb,array,longtable,booktabs,xcolor,hyperref}
\definecolor{reportblue}{HTML}{174BA8}
\hypersetup{colorlinks=true,linkcolor=reportblue,urlcolor=reportblue,pdftitle={PartAware-SG 综合研发与实验报告}}
\setlength{\emergencystretch}{3em}
\setlength{\tabcolsep}{3pt}
\renewcommand{\arraystretch}{1.25}
\title{PartAware-SG 综合研发与实验报告}
\author{}
\date{2026-10-04}
\begin{document}
\maketitle
\tableofcontents
\clearpage

更新：2026-10-04。本文合并原实验记录和集成记录，保留与当前 ScanNet/Hypersim 主线有关的结论。第 4 节依据实际代码解释局部观测与部件融合；本次核对论文及官方评测源码，修订第 6、9 节并新增第 10 节，确定已有评测协议及研究优先级，继续提供离线 LaTeX 渲染版。本次没有生成新的精度结果。完整旧记录及旧运行产物已移到 \texttt{datasets/\allowbreak{}scannet-sg-processed/\allowbreak{}partaware\_\allowbreak{}v1/\allowbreak{}history}。



\textbf{公式阅读：} 打开 \href{RESEARCH\_LOG.html}{LaTeX 渲染版报告} 查看已排版的公式。该 HTML 内嵌由 MathJax 从本报告 LaTeX 生成的矢量公式，可离线阅读；本 Markdown 保留可编辑的 LaTeX 源码。



\section*{1. 目标、备份与目录}
\addcontentsline{toc}{section}{1. 目标、备份与目录}

目标是提高实例与部件建图的可靠性，区分独立物体、真实部件、同一物体的局部观测；尤其避免把两台功放并入桌子。新增部件层与旧物体图共存，基础 ScanNet 建图和 Hypersim 清单接口保留，默认物体关联仍为 legacy。



\begin{itemize}
\item 核心项目：\texttt{/\allowbreak{}home/\allowbreak{}lewisliu/\allowbreak{}PartAware-SG}。


\item 改造前全量备份：\texttt{/\allowbreak{}home/\allowbreak{}lewisliu/\allowbreak{}ScanNet-SG-backup-20261003-175312.\allowbreak{}tar}，包含隐藏文件、Git 状态、旧修改、权重和构建产物；2217 个条目逐项校验。SHA-256：\texttt{b183483537ace34fa854679fb6d4eb7d206312bae6d955d4bf3fd44bc970e8f6}。


\item 原库已恢复到 GitHub 提交 \texttt{52766bba5b4bf7e4773480d1b6c9703c8d8069da}，受版本管理代码保持原状。


\item 当前仅保留 ScanNet/Hypersim 专属路线；移除 3RScan、KITTI-360 适配器、词表和专项测试文档。其他原始数据目录未改动。


\item 第一版结果统一存放在 \texttt{/\allowbreak{}home/\allowbreak{}lewisliu/\allowbreak{}datasets/\allowbreak{}scannet-sg-processed/\allowbreak{}partaware\_\allowbreak{}v1}，与其他实验同级。项目中不保留本轮运行结果、旧编译对照、代码备份及临时下载文件；必需的模型、依赖和兼容补丁保留。

\end{itemize}


已阅读保存 HTML 中的全部可见对话；页面未加载的更早历史不在文件内。归档位于上述外部 history/records。



\section*{2. 阅读与实际接入}
\addcontentsline{toc}{section}{2. 阅读与实际接入}

{\small
\begin{longtable}{|p{\dimexpr\linewidth/2-2\tabcolsep-2\arrayrulewidth\relax}|p{\dimexpr\linewidth/2-2\tabcolsep-2\arrayrulewidth\relax}|}
\hline
来源 & 接入或核对内容 \\ \hline
\endfirsthead
\hline
来源 & 接入或核对内容 \\ \hline
\endhead
\href{https://arxiv.org/abs/2605.10484}{OpenSGA} / \href{https://github.com/tud-amr/ScanNet-SG}{ScanNet-SG} & 核对标签、DINO 256 维特征、SBERT 384 维特征、投影、物体图格式 \\ \hline
\href{https://arxiv.org/abs/2606.29786}{OP3DSG} / \href{https://github.com/AutoCompSysLab/OP3DSG}{代码} & 内置官方物体—部件知识库，改编几何/语义/颜色关联和物体—部件组织 \\ \hline
\href{https://arxiv.org/abs/2305.11173}{VLPart} / \href{https://github.com/facebookresearch/VLPart}{代码} & 官方 Swin-base LVIS/PACO 模型、动态类别分类器、原生掩码 \\ \hline
\href{https://arxiv.org/abs/2309.16650}{ConceptGraphs} / \href{https://github.com/concept-graphs/concept-graphs}{代码} & 掩码区域特征、方向性几何覆盖，可选 DBSCAN 和包含掩码清理 \\ \hline
\end{longtable}
}

代码版本和模型校验见 \texttt{docs/\allowbreak{}sources.\allowbreak{}lock.\allowbreak{}json}、\texttt{docs/\allowbreak{}model\_\allowbreak{}checksums.\allowbreak{}json}。这里只改编 OP3DSG 的先验建图部分，未接入其 LLM 统一推理，不宣称完整复现。未找到可核实的 g-ch/OpenSGA 公共匹配推理仓库，未创建猜测接口。



旧阶段已有 Qwen/联合打分、Hypersim 输入、地板过滤及重复实例处理；本轮保留这些代码。此前 0.395534 候选与 0.39 阈值实验只属于旧结果，本轮不据此改动默认前端阈值。



\section*{3. 接口与改动}
\addcontentsline{toc}{section}{3. 接口与改动}

入口为 \texttt{scannet/\allowbreak{}script/\allowbreak{}run\_\allowbreak{}partaware.\allowbreak{}py}：\texttt{--image-dir} 读取 ScanNet，\texttt{--manifest} 读取包括 Hypersim 的既有 \texttt{scannet\_\allowbreak{}sg\_\allowbreak{}input} 清单；\texttt{--processed-scene} 提供旧物体图及真实实例映射，\texttt{--output} 指向新的外部结果目录。



输出 \texttt{partaware\_\allowbreak{}graph.\allowbreak{}json} 保留所有旧图字段，附加 \texttt{part\_\allowbreak{}nodes}、\texttt{part\_\allowbreak{}relations} 和来源信息。逐帧布尔掩码、观测点云、参数与中文运行日志分别保存为 NPZ、NPY、\texttt{run\_\allowbreak{}config.\allowbreak{}json}、\texttt{run\_\allowbreak{}zh.\allowbreak{}jsonl}。部件使用独立的 CLIP RN50 1024 维特征，不写进物体的 256/384 维特征；旧 TopologyMap 可正常读取物体字段。



部件层的主要保护：父节点必须有类别及掩码覆盖证据；歧义保持未绑定；父编号或部件标签不同不融合；同帧已入轨的部件不再接收其他观测；默认至少两帧确认后才形成 part\_of。已测 Hypersim 的两台功放保持独立实例。这些部件约束不能视为 legacy 物体融合的全局保证，具体差异见第 4 节。



可选项：SAM 框精修、\texttt{--dbscan-eps} 最大簇去噪、\texttt{--subtract-contained} 同真实父实例下的包含清理、\texttt{--erode-pixels} 边界收缩。附加清理默认关闭。C++ 新增 \texttt{--association\_\allowbreak{}mode legacy|op3dsg}，原位置参数及 manifest/filter\_floor 接口保留。旧 SentenceTransformer 构造参数兼容修复不更换模型或维数。



\section*{4. 数学原理与边界}
\addcontentsline{toc}{section}{4. 数学原理与边界}

本节公式整理自当前实现，不把论文中的完整方法直接当成已实现内容。需要区分两个问题：\textbf{同一物体的局部观测融合}是在不同帧中恢复同一实例的身份；\textbf{部件融合}是在同一父物体下关联同一个真实部件。功放的左半边与右半边可能属于一个物体轨迹；桌腿则应保留部件节点，通过 \texttt{part\_\allowbreak{}of} 连接桌子。类别相同、空间包含或颜色相近，都不足以单独证明二者属于同一物理实例。



\subsection*{4.1 统一到世界坐标：融合成立的前提}
\addcontentsline{toc}{subsection}{4.1 统一到世界坐标：融合成立的前提}

对帧 $t$ 的深度像素 $(u,v)$，深度为光轴 Z，$z=d/s$，$s$ 为深度缩放因子。利用深度内参 $K_d$ 和相机到世界位姿 $T_t=[R_t,t_t]$：



\[
\begin{aligned}
p_c&=zK_d^{-1}[u,v,1]^T\\
&=[(u-c_x)z/f_x,\ (v-c_y)z/f_y,\ z]^T,\\
p_w&=R_tp_c+t_t.
\end{aligned}
\]

彩色/深度使用各自内参对齐掩码采样；部件投影先将深度像素映射到彩色像素，再检查其是否属于掩码。沿用原库已配准 RGB-D 的假设，没有另行估计彩色—深度外参；坐标单位为米，无深度和无观测区域不补造几何。后续比较的是世界坐标中的观测点云 $A$ 与已有轨迹点云 $B$，而非不同相机坐标中的点。



小误差下可用下式解释位姿对融合的影响：



\[
\|\delta p_w\|\lesssim
\|\delta p_c\|+\|p_c\|\,\|\delta\theta\|+\|\delta t\|.
\]

旋转误差随距离放大，例如距离约 2 米、旋转误差约 $1^\circ$ 时，仅旋转项就约为 3.5 厘米。这是误差传播示例，不是本次实测值；它解释了为何位姿矫正要先于融合，也说明不能只靠提高语义相似度修复错位点云。



\subsection*{4.2 局部覆盖：为什么只看到一部分也能匹配}
\addcontentsline{toc}{subsection}{4.2 局部覆盖：为什么只看到一部分也能匹配}

定义点 $a$ 到点云 $B$ 的最近邻距离与方向覆盖率：



\[
\begin{aligned}
d(a,B)&=\min_{b\in B}\|a-b\|_2,\\
c_r(A\to B)&=\frac{1}{|A|}\sum_{a\in A}\mathbf1[d(a,B)\le r].
\end{aligned}
\]

记 $\alpha=c_r(A\to B)$、$\beta=c_r(B\to A)$，当前代码有两种几何量：



\[
\begin{aligned}
H&=\begin{cases}\dfrac{2\alpha\beta}{\alpha+\beta},&\alpha+\beta>0,\\0,&\alpha+\beta=0,\end{cases}\\
G&=\max(\alpha,\beta).
\end{aligned}
\]

$H$ 是双向邻域覆盖率的调和平均，任一方向低都会压低得分；$G$ 允许“小观测被大轨迹覆盖”时仍保留高得分。C++ 函数虽名为 \texttt{compute3DMIoU}，返回的首项实际是 $H$，\textbf{不是体积或体素集合的交并比}。邻域覆盖也不是点的一一对应：多个点可以由同一个近邻支持，数值会受到采样密度、掩码污染和位姿误差影响。



例如，小观测完全落在大轨迹上，设 $\alpha=1,\beta=0.1$，则 $H\approx0.182$、$G=1$。这说明较严格的对称覆盖会把同一物体的局部观测拆成新轨迹，最大方向覆盖有助于接回该观测；但桌边与相邻物体的局部接触也可能产生高 $G$，所以仍需要语义与实例约束。



\textbf{当前半径不同：物体 C++ 关联使用 $r=0.10$ 米，部件 Python 关联默认使用 $r=0.03$ 米。} 3 厘米用于细部件，10 厘米用于基础物体融合；二者不能混写为同一参数。



\subsection*{4.3 基础物体的逐帧关联：legacy 与 op3dsg}
\addcontentsline{toc}{subsection}{4.3 基础物体的逐帧关联：legacy 与 op3dsg}

物体关联的语义量来自 \textbf{384 维 SBERT 文本特征} $e$，不是 256 维 DINO 特征。C++ 使用：



\[
q_o=\frac{e_A^Te_B}{\|e_A\|_2\|e_B\|_2+10^{-6}}.
\]

对于一个有效观测，先筛选候选，再选择得分最高的轨迹；无候选则创建新全局编号。当前模式如下，所有条件均需同时成立：



{\small
\begin{longtable}{|p{\dimexpr\linewidth/4-2\tabcolsep-2\arrayrulewidth\relax}|p{\dimexpr\linewidth/4-2\tabcolsep-2\arrayrulewidth\relax}|p{\dimexpr\linewidth/4-2\tabcolsep-2\arrayrulewidth\relax}|p{\dimexpr\linewidth/4-2\tabcolsep-2\arrayrulewidth\relax}|}
\hline
物体模式 & 候选条件 & 候选排序分数 & 同帧排斥 \\ \hline
\endfirsthead
\hline
物体模式 & 候选条件 & 候选排序分数 & 同帧排斥 \\ \hline
\endhead
\texttt{legacy}（默认） & $q_o\ge0.8,\ H\ge0.3$ & $Q_o=0.5q_o+0.5H$ & 逐帧匹配未加入该条件 \\ \hline
\texttt{op3dsg}（可选） & $q_o\ge0.8,\ G\ge0.2,\ Q_o\ge1.2$ & $Q_o=G+(q_o+1)/2$ & 要求当前帧不在候选轨迹的已观测帧集合中 \\ \hline
\end{longtable}
}

物体候选没有额外的标签字符串完全相同条件，可能关联文字嵌入相近但标签不同的观测。\texttt{op3dsg} 物体分支没有使用下文的颜色项；其“方向覆盖 + 语义”与部件的三线索评分是两套实现。两种模式都是按处理顺序进行的贪心关联，没有全局一对一优化。



默认 Florence 改善的是前端裁剪区域的类别/描述证据，随后仍进入上述 C++ 关联；它没有直接替代几何门控或证明物体身份。两个功放即使描述相同，也不能仅凭描述被认定为同一个。



\subsection*{4.4 物体轨迹的最终合并：当前判据及保护范围}
\addcontentsline{toc}{subsection}{4.4 物体轨迹的最终合并：当前判据及保护范围}

逐帧处理后，C++ 还会合并已有物体轨迹。先以预先计算的点云中心距离 $\le2$ 米筛选，再计算同样的 $q_o,\alpha,\beta,H$。设 $F_A,F_B$ 为不同观测帧的集合，$l=\min(\alpha,\beta)$、$h=\max(\alpha,\beta)$。



当前 \texttt{legacy} 的原合并条件为：



\[
M_0=(H>0.7)\lor
\big[q_o>0.7\ \land\ (H>0.3\ \lor\ \alpha>0.8\ \lor\ \beta>0.8)\big].
\]

另外有几何恢复条件，尝试找回被中等文本相似度拆开的重叠轨迹：



\[
\begin{aligned}
M_{rec}={}&\big[F_A\cap F_B=\varnothing\big]\\
&\land\big[|F_A|\ge2,\ |F_B|\ge2\big]\\
&\land\big[q_o\ge0.50,\ l\ge0.50,\ h\ge0.80\big].
\end{aligned}
\]

\texttt{legacy} 最终允许 $M_0\lor M_{rec}$。只有恢复分支要求无共享帧，$M_0$ 不要求；其 $H>0.7$ 分支甚至不要求语义阈值。因此不能宣称默认物体算法能够严格防止同帧物体合并，或每次合并都经过强语义验证。



可选 \texttt{op3dsg} 的轨迹合并使用另一判据：



\[
\begin{aligned}
M_{op}={}&\big[F_A\cap F_B=\varnothing\big]\\
&\land\big[q_o\ge0.8,\ h\ge0.7,\ l\ge0.1\big].
\end{aligned}
\]

该模式不使用 $M_0\lor M_{rec}$。共享帧排斥较保守：既能保护同帧独立物体，也会阻止同一物体在同帧被误分出的两块直接合并。最终编号沿 \texttt{merge\_\allowbreak{}map} 链重映射，少于 2 个不同帧支持的物体轨迹被移除。中心在合并循环前只计算一次，合并按遍历顺序进行；当前没有合并撤销和全局最优保证。



\subsection*{4.5 部件融合：几何、CLIP 语义与颜色}
\addcontentsline{toc}{subsection}{4.5 部件融合：几何、CLIP 语义与颜色}

部件观测先按 1 厘米体素压缩点云。其候选集合要求\textbf{部件标签相同、父编号相同、当前帧尚未进入候选轨迹}，然后计算：



\[
\begin{aligned}
\hat f&=\frac{f}{\|f\|_2},\\
q_p&=\hat f_A^T\hat f_B,\\
S_p&=\frac{q_p+1}{2}.
\end{aligned}
\]

$f$ 是独立的 CLIP RN50 掩码区域 1024 维特征，不能与物体的 DINO/SBERT 特征混用。几何使用 $G$，硬门控为 $G\ge0.2$ 且 $q_p\ge0.7$。



颜色用四个通道的 32 桶归一化直方图。将 RGB 缩放至 $[0,1]$，颜色通道为：



\[
\begin{aligned}
c_1&=\operatorname{clip}\!\left(\frac{a_{Lab}+128}{256},0,1\right),\\
c_2&=\operatorname{clip}\!\left(\frac{b_{Lab}+128}{256},0,1\right),\\
c_3&=\frac{R-G+1}{2},\\
c_4&=\frac{R+G-2B+2}{4}.
\end{aligned}
\]

每通道直方图 $h^{(c)}$ 满足 $\sum_kh_k^{(c)}=1$。对等间距桶中心 $x_k=(k-0.5)/32$，一维 Wasserstein 距离可写为累积分布差的面积：



\[
\begin{aligned}
F_A^{(c)}(k)&=\sum_{i=1}^{k}h_{A,i}^{(c)},\qquad
F_B^{(c)}(k)=\sum_{i=1}^{k}h_{B,i}^{(c)},\\
W_1(h_A^{(c)},h_B^{(c)})
&=\frac{1}{32}\sum_{k=1}^{31}\left|F_A^{(c)}(k)-F_B^{(c)}(k)\right|,\\
D&=\frac14\sum_{c=1}^{4}W_1(h_A^{(c)},h_B^{(c)}),\\
C&=\frac{1}{1+3D}.
\end{aligned}
\]

这里的 $F_A^{(c)}$、$F_B^{(c)}$ 是颜色直方图的累积分布；与第 4.4 节中的观测帧集合区分。



该距离衡量颜色分布移动需要的代价，邻近色桶的变化比相隔很远的色桶变化代价小。它使用 LAB/opponent 通道，不等同于论文完整的 Color Names 编码；光照、反光和背景混入仍会改变分布。



通过硬门控后，当前部件分数为：



\[
\begin{aligned}
Q_p&=2[(1-w)(G+S_p)+wC],\\
w&=0.6,\qquad Q_p\ge1.5.
\end{aligned}
\]

在默认权重下等价于 $Q_p=0.8G+0.8S_p+1.2C$。分数最高且达到门限的候选被更新，否则新建部件轨迹。前面的因子 2 是实现中的评分尺度，分数最高可达 2.8，\textbf{不是概率，也不是 $[0,1]$ 的置信度}。颜色无法绕过几何/语义门控，检测置信度也不直接参与该关联分数；它单独保存在轨迹中。



\subsection*{4.6 父节点归属：覆盖证据不等于身份合并}
\addcontentsline{toc}{subsection}{4.6 父节点归属：覆盖证据不等于身份合并}

对部件掩码 $M_p$，仅考虑名称与其父类别一致、且编号存在于基础物体图的候选父实例。父实例 $j$ 的局部掩码记为 $M_j$：



\[
a_j=\frac{|M_p\cap M_j|}{|M_p|},\qquad
j^*=\arg\max_j a_j.
\]

同一全局父编号对应多个局部候选时，先取其最大覆盖率再排序。只有 $a_{j^*}\ge0.2$，且存在第二候选时满足 $a_{j^*}-a_{j^{(2)}}\ge0.1$，才绑定父节点；否则 \texttt{parent\_\allowbreak{}id=None}。只有一个候选时不做第二名差值判断。这是部件像素被父掩码覆盖的比例，不是二维 IoU，也不是经过校准的归属概率。



部件关联要求父编号严格相等，避免相邻桌子的桌腿融合。但两个未绑定轨迹的父编号都为 \texttt{None} 时，也满足“父编号相等”，仍可在标签、帧和相似度门控下融合；当前没有后续重新绑定父节点的步骤。无父轨迹与已绑定轨迹不能直接融合，单帧归属变化也可能导致碎片，这正是软父归属的后续研究动机。



轨迹确认和关系导出分别为：



\[
\begin{aligned}
\mathrm{confirmed}(k)&\iff |F_k|\ge2,\\
\mathrm{part\_of}(k,j)&\iff \mathrm{confirmed}(k)\\
&\qquad\land(\mathrm{parent\_id}_k=j\ne\mathrm{None}).
\end{aligned}
\]

确认采用不同帧数，不采用点数或单帧重复检测数；相邻重复视角和持续误检仍可能通过该条件，不能把 \texttt{confirmed} 当成真值。关系导出增加子节点和父子边，不把部件的 1024 维特征覆盖进基础物体接口。



\subsection*{4.7 合并后的点云与特征如何更新}
\addcontentsline{toc}{subsection}{4.7 合并后的点云与特征如何更新}

\textbf{部件点云。} 设体素大小 $v=0.01$ 米，点的体素索引为 $\kappa(p)=\lfloor p/v\rfloor$。每个非空体素用其输入点的质心表示：



\[
\begin{aligned}
P_k&=\{p\in P:\kappa(p)=k\},\\
\mu_k&=\frac{1}{|P_k|}\sum_{p\in P_k}p,\\
\mathcal V_v(P)&=\{\mu_k:|P_k|>0\},\\
B_{new}&=\mathcal V_v(B_{old}\mathbin{\Vert}\mathcal V_v(A)).
\end{aligned}
\]

$\Vert$ 表示拼接。压缩后比较邻域覆盖可降低点密度差异及计算量，但实现未保存体素历史点数：旧质心与新输入质心再次等权参与计算，不等价于对全部原始点一次性求质心，也不是 TSDF 或位姿优化。



\textbf{部件特征与置信度。} 已有 $n$ 次观测、当前特征为 $\hat f_n$，新归一化观测为 $\hat f_{obs}$：



\[
\begin{aligned}
\hat f_{n+1}&=\frac{n\hat f_n+\hat f_{obs}}{\|n\hat f_n+\hat f_{obs}\|_2},\\
c_{n+1}&=\frac{nc_n+c_{obs}}{n+1}.
\end{aligned}
\]

前者是\textbf{归一化递推更新}，不是严格的“全部原始特征求和后再归一化”，因为历史向量每步已被重新归一化，代码没有保存原始特征和。后者是检测置信度的算术平均，但检测分数本身不保证概率校准。颜色逐通道采用指数平滑并重新归一化：



\[
\begin{aligned}
\tilde h_{n+1}^{(c)}&=0.7h_n^{(c)}+0.3h_{obs}^{(c)},\\
h_{n+1}^{(c)}&=\frac{\tilde h_{n+1}^{(c)}}{\sum_k\tilde h_{n+1,k}^{(c)}}.
\end{aligned}
\]

因此新颜色观测总占 0.3，历史单次观测的贡献随更新衰减，不是按全部帧均匀平均。



\textbf{物体点云与代表特征。} C++ 物体点云直接拼接，逐帧更新后的点数超过 10000 时按索引隔一个取一个，并非采用上述部件体素操作。在线 DINO 特征使用计数均值：$\bar f_{n+1}=(n\bar f_n+f_{obs})/(n+1)$。文本嵌入、名称、描述和置信度则来自“单次观测点数最多”的代表观测，未做 SBERT 均值或置信度平均；点数多只是一项启发式，不能保证观测质量最好。



还需记录一项现有实现限制：最终合并轨迹 $A,B$ 时，DINO 更新实际为



\[
\begin{aligned}
\bar f_{merge}^{current}&=\frac{n_A\bar f_A+\bar f_B}{n_A+1},\\
n_{new}&=n_A+n_B.
\end{aligned}
\]

它把另一条轨迹的均值当作一次观测，当 $n_B>1$ 时不等于双方全部观测的加权均值。若要保持计数平均，应使用 $(n_A\bar f_A+n_B\bar f_B)/(n_A+n_B)$；这是\textbf{待修正项，本次只修订报告，未修改融合代码}。同样，当前递推特征和体素代表点会受处理顺序影响，报告不把它们描述为无偏估计或全局最优融合。



\subsection*{4.8 数值示例、失效条件与预期收益}
\addcontentsline{toc}{subsection}{4.8 数值示例、失效条件与预期收益}

以下是假设值，用来解释判据，不是新增实验结果。采用第 4.2 节的局部观测 $\alpha=1,\beta=0.1$，设 $q_o=q_p=0.9$、$C=0.8$，并假定所有标签、父编号及同帧条件满足：



{\small
\begin{longtable}{|p{\dimexpr\linewidth/3-2\tabcolsep-2\arrayrulewidth\relax}|p{\dimexpr\linewidth/3-2\tabcolsep-2\arrayrulewidth\relax}|p{\dimexpr\linewidth/3-2\tabcolsep-2\arrayrulewidth\relax}|}
\hline
关联路线 & 代入结果 & 解释 \\ \hline
\endfirsthead
\hline
关联路线 & 代入结果 & 解释 \\ \hline
\endhead
legacy 物体逐帧关联 & $H\approx0.182<0.3$ & 本次在线匹配失败；后续轨迹合并仍需另按第 4.4 节判断 \\ \hline
op3dsg 物体逐帧关联 & $G=1,\ Q_o=1+0.95=1.95$ & 门控满足时可以接回局部观测 \\ \hline
部件关联 & $Q_p=0.8+0.76+0.96=2.52$ & 达到 1.5 门限，同一部件可累积跨帧几何 \\ \hline
\end{longtable}
}

相反，即使两个功放的语义和颜色高度相似，若 $G=0$，上述在线几何门控仍不能通过。若同一物体前后两侧完全没有共同可见点，$G$ 也可能接近 0，现有方法仍会分裂；若掩码混入桌面且位姿错位，局部高覆盖又可能造成误合并。因此方向覆盖的预期收益是减少\textbf{有几何重叠的局部观测碎片}，并非保证所有部分可见物体都能恢复完整身份。



可选清理仍独立于身份判定：包含比例达到 0.9 且面积比超过 1.5 时，只对同已解析父实例、同父类别且不同部件标签的掩码扣除包含区域。DBSCAN 保留最大簇可能清理背景，也可能删除真实不连通结构；两项默认关闭，需分别消融。



评价应区分错误拆分与错误合并：前者把一个真实物体分成多轨迹，后者把多个真实物体合成一轨迹；同时检查父归属正确率、部件召回与几何误差。减少节点数不能单独证明建图精度提高；此前实验未完成真值评测，不新增精度或 SOTA 声明。多视角 Florence、软父归属和可撤销合并仍属于第 9 节的规划，评测协议见第 10 节。



\subsection*{4.9 公式与实现的对应位置}
\addcontentsline{toc}{subsection}{4.9 公式与实现的对应位置}

本次以项目提交 \texttt{c7febf0} 下的当前源文件逐项核对：



{\small
\begin{longtable}{|p{\dimexpr\linewidth/2-2\tabcolsep-2\arrayrulewidth\relax}|p{\dimexpr\linewidth/2-2\tabcolsep-2\arrayrulewidth\relax}|}
\hline
实现 & 本节对应内容 \\ \hline
\endfirsthead
\hline
实现 & 本节对应内容 \\ \hline
\endhead
\texttt{scannet/\allowbreak{}src/\allowbreak{}openset\_\allowbreak{}ply\_\allowbreak{}map.\allowbreak{}cpp}：\texttt{cosineSimilarity}、\texttt{compute3DMIoU}、\texttt{matchInstanceToGlobal}、最终合并循环 & 10 厘米物体几何、SBERT 关联、两种模式的在线/离线门控、编号重映射及 DINO 更新限制 \\ \hline
\texttt{scannet/\allowbreak{}script/\allowbreak{}partaware/\allowbreak{}fusion.\allowbreak{}py}：\texttt{overlap}、\texttt{color\_\allowbreak{}histogram}、\texttt{color\_\allowbreak{}similarity}、\texttt{PartFusion.\allowbreak{}add/\allowbreak{}export} & 3 厘米部件关联、CLIP/颜色评分、递推更新及确认/父子边导出 \\ \hline
\texttt{scannet/\allowbreak{}script/\allowbreak{}partaware/\allowbreak{}geometry.\allowbreak{}py}：\texttt{project\_\allowbreak{}mask}、\texttt{voxel\_\allowbreak{}downsample} & 世界坐标投影与体素质心 \\ \hline
\texttt{scannet/\allowbreak{}script/\allowbreak{}run\_\allowbreak{}partaware.\allowbreak{}py}：\texttt{select\_\allowbreak{}parent} & 掩码覆盖、全局父编号去重和归属歧义判据 \\ \hline
\end{longtable}
}

本次验证范围为报告与源码的公式、阈值及分支一致性，没有重跑模型或生成新场景图；后续改动默认参数时，也应同步更新本节。



\section*{5. 数据集实测}
\addcontentsline{toc}{section}{5. 数据集实测}

原始输入固定来自 \texttt{/\allowbreak{}home/\allowbreak{}lewisliu/\allowbreak{}datasets}。Hypersim 使用用户指定的 \texttt{scannet-sg-input/\allowbreak{}hypersim/\allowbreak{}ai\_\allowbreak{}001\_\allowbreak{}002}，100 帧中 98 帧有完整旧实例映射；\texttt{cam\_\allowbreak{}02\_\allowbreak{}0024}、\texttt{cam\_\allowbreak{}03\_\allowbreak{}0034} 缺少旧 updated\_instance 文件，明确记录并跳过。ScanNet 使用真实 \texttt{scene0000\_\allowbreak{}00} 的 30 帧。



{\small
\begin{longtable}{|p{\dimexpr\linewidth/6-2\tabcolsep-2\arrayrulewidth\relax}|p{\dimexpr\linewidth/6-2\tabcolsep-2\arrayrulewidth\relax}|p{\dimexpr\linewidth/6-2\tabcolsep-2\arrayrulewidth\relax}|p{\dimexpr\linewidth/6-2\tabcolsep-2\arrayrulewidth\relax}|p{\dimexpr\linewidth/6-2\tabcolsep-2\arrayrulewidth\relax}|p{\dimexpr\linewidth/6-2\tabcolsep-2\arrayrulewidth\relax}|}
\hline
分支 & 帧数 & 部件轨迹 & 多帧确认 & part\_of & 耗时（秒） \\ \hline
\endfirsthead
\hline
分支 & 帧数 & 部件轨迹 & 多帧确认 & part\_of & 耗时（秒） \\ \hline
\endhead
hypersim：VLPart 原掩码 & 98 & 30 & 20 & 17 & 20.09 \\ \hline
hypersim：SAM+去噪+包含清理 & 98 & 27 & 18 & 15 & 109.83 \\ \hline
scannet：VLPart 原掩码 & 30 & 14 & 9 & 9 & 13.85 \\ \hline
\end{longtable}
}

C++ 扩大样本回归：Hypersim 旧二进制/新 legacy 均为 11 个物体节点，op3dsg 为 14 个；ScanNet 三者均为 4 个。两种数据的新 legacy 完整 JSON 都与改造前保留的旧二进制精确一致。这里的重新融合图与用于部件实验的旧清理后图属于不同阶段，节点数不直接混比。



v1 的 15 项单元测试、shell 语法及真实图接口检查通过。所有新部件图的旧字段逐项保持，原读取器可读；Hypersim 中两台功放 id=3、11 保持独立。之前原 ScanNet 分割一帧也实测通过，输出仍为 256/384 维。



上述 v1 图、逐帧结果、中文运行日志、机器可读摘要和测试输出位于外部 \texttt{partaware\_\allowbreak{}v1}。模型环境使用隔离的 \texttt{.\allowbreak{}venv-vlpart}，避免改动原 conda 的 timm；具体安装见 \texttt{docs/\allowbreak{}PARTAWARE\_\allowbreak{}SETUP.\allowbreak{}md}。



\textbf{补充已完成的 scene0802\_00 全帧实验。} 输入为用户 ZED 数据经提供的 COLMAP 脚本矫正后保留的 443 帧，不是带官方真值标注的 ScanNet 场景。两套结果均使用该有效帧集和 legacy 物体关联，分别保存于外部 \texttt{scene0802\_\allowbreak{}scannetsg\_\allowbreak{}v1}、\texttt{scene0802\_\allowbreak{}partaware\_\allowbreak{}v1}：



{\small
\begin{longtable}{|p{\dimexpr\linewidth/5-2\tabcolsep-2\arrayrulewidth\relax}|p{\dimexpr\linewidth/5-2\tabcolsep-2\arrayrulewidth\relax}|p{\dimexpr\linewidth/5-2\tabcolsep-2\arrayrulewidth\relax}|p{\dimexpr\linewidth/5-2\tabcolsep-2\arrayrulewidth\relax}|p{\dimexpr\linewidth/5-2\tabcolsep-2\arrayrulewidth\relax}|}
\hline
路线 & 物体节点 / 关系边 & 部件轨迹 & 不同帧确认 & 有父节点的确认部件 / part\_of \\ \hline
\endfirsthead
\hline
路线 & 物体节点 / 关系边 & 部件轨迹 & 不同帧确认 & 有父节点的确认部件 / part\_of \\ \hline
\endhead
DINO/SAM 基础入口 & 21 / 85 & — & — & — \\ \hline
默认 Florence + VLPart/SAM 部件层 & 31 / 211 & 580 & 386 & 197 \\ \hline
\end{longtable}
}

PartAware 还有 189 条已确认但无父归属的轨迹、194 条单帧轨迹，默认可视化不显示它们。全帧掩码、特征维数和部件图中保留的原物体字段已核验；详细运行记录见两个实验根目录的 \texttt{实验报告.\allowbreak{}md}。该对比同时改变物体前端和部件层，不能用节点数差异单独证明融合算法的精度收益；本次报告修订未重新运行实验。



\section*{6. 结论与研究高度}
\addcontentsline{toc}{section}{6. 结论与研究高度}

交付已具备可运行、可消融、兼容旧接口的部件建图基线。当前是模型与已有方法的集成，尚不能称为独立算法贡献或建图精度的本质突破。Florence 增强候选语义证据，VLPart 增加部件感知，OP3DSG 提供知识库和融合设计参考；默认物体关联仍为 legacy，附加部件层不回头纠正已有物体身份。此前轨迹数量与运行成功只能证明功能可用。



本次发现本地 ScanNet 实例真值，但尚未完成精度评测；部件、父子及功能关系真值仍需补齐。下一阶段采用第 10 节已有协议，固定前端和帧集后分别消融关联、颜色、SAM、去噪与父归属。潜在研究问题是：在部分可见、同类物体相邻和粒度不一致时，如何联合维护实例身份与部件归属，并在新证据到来后纠错。第 9 节将其分解为可验证步骤；这些步骤本身已有相关研究，独立贡献仍需新的方法差异、强基线与多场景证据。



本轮新增代码、注释及 README 扩展使用英语，综合报告和运行消息使用中文。



\section*{7. 共用原版可视化（v2 修正）}
\addcontentsline{toc}{section}{7. 共用原版可视化（v2 修正）}

v1 曾使用独立的简化部件绘图，没有完整复用原版的颜色、包围盒和交互；该实现已经替换。现在正常图和精修图均调用 \texttt{script/\allowbreak{}utils/\allowbreak{}result\_\allowbreak{}visualization.\allowbreak{}py:\allowbreak{}:\allowbreak{}visualize\_\allowbreak{}map\_\allowbreak{}with\_\allowbreak{}nodes}，只通过 \texttt{build\_\allowbreak{}part\_\allowbreak{}overlay} 追加子节点和虚线父子边。



不开部件时，真实 Hypersim 数据的 41 个基础几何对象（含点云/颜色/节点/包围盒/关系）与提交 adaa808 中原绘图函数的数组及哈希逐项一致；开部件时这 41 个对象仍逐项一致，只追加 30 个几何：15 个已确认且有父实例的子节点与 15 条父子边。小节点采用父实例配色，物体关系保持原色，默认隐藏部件点云与无归属轨迹。诊断开关仍能显示全部 27 条部件轨迹及其点云。



两个入口均复用原 GUI/legacy 窗口，GUI 拾取包含物体与子节点。原版 Open3D 实际生成了截图并作图像检查；拾取回调已接入，未模拟真实鼠标点击。新增可选尾参数和 CLI 见接口文档，旧参数保留；\texttt{visualize\_\allowbreak{}partaware.\allowbreak{}py} 是兼容包装器。20 项部件与显示回归检查通过；旧 Florence 参考配置及显式 DINO 回退也通过真实图像测试。



所有可视化命令均带用户要求的 \texttt{env -u WAYLAND\_\allowbreak{}DISPLAY -u WAYLAND\_\allowbreak{}SOCKET} 和 \texttt{XDG\_\allowbreak{}SESSION\_\allowbreak{}TYPE=x11 LIBGL\_\allowbreak{}ALWAYS\_\allowbreak{}SOFTWARE=true}。使用 tracking-ID 编码的 \texttt{instance\_\allowbreak{}cloud*.\allowbreak{}ply}，避免把已经着色的 RGB 导出重新当实例编号解释。



\section*{8. Florence 默认入口与真实 v2 测试}
\addcontentsline{toc}{section}{8. Florence 默认入口与真实 v2 测试}

基本物体流程默认采用 DINO 候选框 → Florence 裁剪区域描述证据 → SAM → 原 C++ 融合/JSON。ScanNet 文件夹入口及 Hypersim manifest 不变，256 维视觉特征与 384 维文字特征不变。\texttt{--grounding\_\allowbreak{}backend dino} 或 runner 的 \texttt{GROUNDING\_\allowbreak{}BACKEND=dino} 可复现原前端，显式 \texttt{--joint\_\allowbreak{}config} 保留旧场景参考归一化与硬负样本逻辑。部件检测仍使用 VLPart/CLIP，并未声称已经由 Florence 替代。



便携默认不依赖本场景参考图；全部输入类别获得描述证据。只实例化一个长期 Florence worker，在已有 \texttt{sg-florence} 环境中运行（transformers 4.49.0 / timm 1.0.15 / einops 0.8.1），避免升级基础 RAM 或 VLPart 依赖。DINO、Florence、SAM 分时进入 GPU。修复了 transformers 把 Florence 条件导入 FlashAttention 当成必需依赖的问题，仅在本模型加载期间使用 eager attention 兼容处理；显式加载可信本地配置消除交互提示。缺失模型或推理错误会明确失败，不悄悄切换后端。



对候选裁剪区域 $I$，描述 $d=(y_1,\ldots,y_N)$ 的证据为普通 token 的平均条件对数似然：



\[
\begin{aligned}
\ell(d\mid I)&=\frac{1}{N}\sum_k\log p(y_k\mid y_{<k},I,\texttt{<CAPTION>}),\\
m&=\ell(d_c\mid I)-\max_{j\ne c}\ell(d_j\mid I).
\end{aligned}
\]

描述只作为 decoder 目标，encoder 不看到目标答案。默认融合为 $s=\sigma(0.8\operatorname{logit}(p_{DINO})+0.2m)$。这只是可解释的未校准证据，不能作为准确概率。默认比较背景及少量易混类别，采用软证据；旧显式参考配置仍按其硬负样本规则运行。



真实测试先发现：短模板配合硬语义否决会误拒绝可见功放。该失败消融已保存在 v2/history；随后改为描述结构特征及软证据。真实 Hypersim 5 帧与 ScanNet 2 帧默认推理完成，并通过原 C++ 融合、generate\_json、SAM 部件精修完成全链路。新物体图分别为 9 / 3 个节点，部件轨迹为 13 / 1 条；其旧图字段逐项一致。小样本及抽帧跨度会改变观测支持，不能与 v1 的 98 / 30 帧节点数直接比较。



v2 输出为 \texttt{datasets/\allowbreak{}scannet-sg-processed/\allowbreak{}partaware\_\allowbreak{}v2/\allowbreak{}\{hypersim/\allowbreak{}ai\_\allowbreak{}001\_\allowbreak{}002,scannet/\allowbreak{}scene0000\_\allowbreak{}00\}}，与其他实验同级。这是 7 帧端到端功能验证，未重跑完整 100 帧；v1 全帧结果和原始失败记录保留。新代码修改不会自动更新旧缓存。



功放诊断：v1 id=3、11 的中心相距约 0.58 米，原始画面有两台物理功放，且两帧 updated\_instance 记录同时含这两个编号。左侧首帧仅有一个功放实例掩码，但边缘未完整覆盖，融合点云还含离散点。因此不能把“两个 amplifier 类别节点”直接判为重复实例；用户看到的左侧碎片需结合具体视角检查掩码缺失、类别竞争与跨帧投影误差。诊断数据和回投影保存在 v2，未强制合并两个节点，也未宣称已经消除所有碎片。



\section*{9. 论文含义与算法优先级（研究方案，尚未实现）}
\addcontentsline{toc}{section}{9. 论文含义与算法优先级（研究方案，尚未实现）}

\subsection*{9.1 三种接入各自解决什么}
\addcontentsline{toc}{subsection}{9.1 三种接入各自解决什么}

{\small
\begin{longtable}{|p{\dimexpr\linewidth/3-2\tabcolsep-2\arrayrulewidth\relax}|p{\dimexpr\linewidth/3-2\tabcolsep-2\arrayrulewidth\relax}|p{\dimexpr\linewidth/3-2\tabcolsep-2\arrayrulewidth\relax}|}
\hline
模块 & 原研究与当前作用 & 对建图的边界 \\ \hline
\endfirsthead
\hline
模块 & 原研究与当前作用 & 对建图的边界 \\ \hline
\endhead
Florence & 当前对 DINO 裁剪候选计算描述证据，辅助类别竞争 & 同类物体可以获得相同描述；描述不能证明实例身份，不能修复位姿或补出不可见表面 \\ \hline
\href{https://openaccess.thecvf.com/content/ICCV2023/html/Sun\_Going\_Denser\_with\_Open-Vocabulary\_Part\_Segmentation\_ICCV\_2023\_paper.html}{VLPart，ICCV 2023} & 研究二维开放词汇部件分割；联合部件、物体及图像级监督，并利用稠密语义对应迁移到新类别。当前使用其预训练模型与动态分类器 & 有助于找把手、按钮等小部件；其 AP/mAP 是二维检测分割指标，不是三维图精度，模型也不负责跨帧身份或关系推理 \\ \hline
\href{https://arxiv.org/html/2606.29786v1}{OP3DSG} & 研究物体、部件、空间/功能关系和可供性统一建图；结合知识引导、几何/语义/颜色融合与受几何约束的 LLM 推理 & 当前仅改编知识库和部分融合先验，未复现完整推理。部件关联颜色度量也不同于论文，不能直接沿用其实验成绩 \\ \hline
\end{longtable}
}

改进潜力来自新增的可观察证据及约束，而不是模型名字或数量。漏检小部件可能由 VLPart 改善；局部观测碎片可能由关联改善；错误功能边需视觉证据和约束校验。这三类误差应分开消融。



\subsection*{9.2 多视角关联：最高优先级，但不是已有流程之外的新概念}
\addcontentsline{toc}{subsection}{9.2 多视角关联：最高优先级，但不是已有流程之外的新概念}

原 ScanNet-SG 已执行 RGB-D 投影、跨帧关联及最终轨迹合并；当前主要限制是贪心决策、覆盖与文字门控、代表特征选择，以及缺少合并撤销。建议先修正第 4.7 节的轨迹均值权重，再实现以下两级关联：



\begin{enumerate}
\item \textbf{帧内先处理分割粒度，再进行跨帧身份匹配。} 完整且互相分离的同帧实例应提供排斥证据；同一物体被切成两块的候选需要连接、轮廓和多视角证据，不能把“同帧禁止融合”当作绝对规则，也不能放开后一律按类别合并。


\item \textbf{跨帧使用可见性与投影一致性，保留拒绝匹配。} 仅将深度检查后可见的历史点回投当前图像，结合掩码覆盖、位置、外观和语义形成稀疏候选。在候选图上做带未匹配选项的一对一分配，减轻贪心顺序影响。完全无公共表面的正反面仍可能无法关联，需要中间视角或额外结构证据。

\end{enumerate}


\href{https://arxiv.org/html/2605.15753v1}{层级功能图论文} 已使用投影、几何、颜色、语义及最大权一对一匹配，因此单独加入 Hungarian 算法不是创新。我们的额外问题是粒度冲突和错误决策可修正性，需要保留观测来源及编号映射，支持撤销错误合并，最后仍导出原接口。



Florence 多视角证据应只在实例身份初步可靠后汇总，选择视角有差异且可见度较好的裁剪；否则误并两台物体后得到“稳定描述”，会把错误强化。相邻重复帧并非独立证据。参考 \href{https://arxiv.org/html/2403.17846v2}{HOV-SG} 的特征去异常思想，稳健聚合与融合规则需分别消融。



\subsection*{9.3 层级优化：对父归属有用，不能代替物体融合}
\addcontentsline{toc}{subsection}{9.3 层级优化：对父归属有用，不能代替物体融合}

当前 \texttt{select\_\allowbreak{}parent} 依靠单帧覆盖并给出固定父编号；后续 \texttt{PartFusion.\allowbreak{}add} 要求父编号严格相同。因此一次归属错误或 \texttt{None} 会使同一部件后续分裂，部件层也不能主动合并两个已有父物体。



建议为部件保留多个父候选及“未知”状态，多帧更新后再确定归属。可参考上述层级论文的历史证据、熵正则和时间平滑，以下为迁移后的设计表达：



\[
\begin{aligned}
L_j^{1:t}&=\sum_{\tau\le t}w_\tau\,\operatorname{logit}(s_{\tau j}),\\
z_t^*&=\arg\max_{z\in\Delta}
\left[z^TL^{1:t}+\lambda_H H(z)
-\frac{\lambda_T}{2}\|z-z_{t-1}\|_2^2\right].
\end{aligned}
\]

这里 $z$ 是父候选分布，熵项保留证据不足时的歧义，平滑项抑制跳变。当前 Florence margin、覆盖率及部件关联分数都不是校准概率，不可直接作为 $s_{\tau j}$；须校准并截断，或另外设计有界证据。候选缺席/遮挡不作为反证，相邻帧支持要折扣，平滑也不能阻止强新证据纠错。



结构约束可检查 \texttt{part\_\allowbreak{}of} 无环、确定后只有一个直接结构父节点、尺度及包含相容。功能关系必须另存，允许多对多与远距离连接：把手属于抽屉不等于控制同场景任意柜子；遥控器和电视的空间接近也不能证明配对。新增功能边只会提高表达能力，只有真值关系指标提高才说明更准确。



\subsection*{9.4 更深入的候选：不确定性融合}
\addcontentsline{toc}{subsection}{9.4 更深入的候选：不确定性融合}

\href{https://arxiv.org/html/2607.07170v1}{PUF} 将关联改为语义与空间似然，并用 Dirichlet 累积节点/关系分布。相较继续叠加检测器，它更直接涉及“如何相信和累积观测”。但\href{https://github.com/yyyyangyi/PUF}{官方实现} 消费完整类别和关系概率；当前开放词汇特征与描述 margin 不具备同一概率接口，不能把其算法直接插入便宣称复现。可研究在固定候选词表下校准证据、保留未知质量、延迟硬决策，并在原字段之外保存不确定性。该方向需要独立消融和身份冲突保护，软关联也可能把证据错误摊到多个相邻实例。



实施顺序：评测基线 → 特征更新修正 → 投影与可见性关联 → 可撤销身份决策 → 软父归属 → 有证据的功能边；概率融合作为单独研究分支。全部保留 ScanNet/Hypersim 输入、旧输出字段与共用可视化。环境版本及模型源码哈希见 \texttt{docs/\allowbreak{}environment\_\allowbreak{}florence.\allowbreak{}json}。



\section*{10. 采用已有评测协议（已选型，尚未运行）}
\addcontentsline{toc}{section}{10. 采用已有评测协议（已选型，尚未运行）}

\subsection*{10.1 来源与选定指标}
\addcontentsline{toc}{subsection}{10.1 来源与选定指标}

已有明确指标，无需另造一个“图看起来更好”的总分。\href{https://arxiv.org/html/2402.12259v2}{Open3DSG，CVPR 2024} 已分别评估物体、谓词和三元组的 R@K、mR@K；开放词汇结果映射到固定词表。其 3DSSG 数据与我们的 ScanNet/Hypersim 不同，复用定义不代表分数可以直接与原论文横向比较，也无需恢复 3RScan 专属接口。



本项目选用下面的分层方案，论文原协议与适配检查分栏报告：



{\small
\begin{longtable}{|p{\dimexpr\linewidth/3-2\tabcolsep-2\arrayrulewidth\relax}|p{\dimexpr\linewidth/3-2\tabcolsep-2\arrayrulewidth\relax}|p{\dimexpr\linewidth/3-2\tabcolsep-2\arrayrulewidth\relax}|}
\hline
评测内容 & 采用指标 & 用途与限制 \\ \hline
\endfirsthead
\hline
评测内容 & 采用指标 & 用途与限制 \\ \hline
\endhead
ScanNet 物体实例 & \href{https://kaldir.vc.in.tum.de/scannet\_benchmark/semantic\_instance\_3d}{官方 AP25、AP50、AP}，AP 对 0.50:0.05:0.95 重叠阈值平均 & 在官方网格顶点上评估实例掩码和类别，重复预测计为假阳性；能揭示拆分、误合并与边界污染。仅覆盖官方 18 个有效实例类别，不把功放等未覆盖类别伪装成该 benchmark 成绩 \\ \hline
开放词汇物体、部件节点 & OP3DSG 节点 R@3、R@5；补报 R@3 的 IoU>0.10、IoP>0.25 定位敏感性 & 保留原论文的宽松 IoU>0 主协议；严格定位列独立报告，IoP 单独不能惩罚过小碎片 \\ \hline
部件归属、功能三元组 & OP3DSG Assoc. Node 和 Triplet R@3、R@5、R@10 & 采用官方标签空间检索定义；无功能预测时不制造可供性或功能分数 \\ \hline
完整有向空间图 & 在固定一对一节点对应下报告三元组 Precision/Recall/F1，按关系类别另报召回 & 属于项目适配评测，检查端点、方向与谓词；完整标注后才统计假阳性。补齐稳定三元组排名后，再按标准 R@50、mR@50 协议评价 \\ \hline
\end{longtable}
}

全图 R@50、mR@50 的精确定义参考 \href{https://arxiv.org/html/2404.09616v1}{场景图指标论文及 SGBench}：从置信排序中取图的前 50 个三元组，并规定每对端点可选几个关系。\textbf{这与 OP3DSG 的标签检索 top-K 不同，禁止混写同名列。} 当前旧关系没有统一可靠排名，不用随意排序来生成 R@50。



\subsection*{10.2 复用时必须保留的细节}
\addcontentsline{toc}{subsection}{10.2 复用时必须保留的细节}

本次检查了 OP3DSG 的\href{https://github.com/AutoCompSysLab/OP3DSG/blob/main/src/eval/eval\_uni.py}{官方评测源码}（访问：2026-10-04）。其标签候选 K 为 1、2、3、5、10；节点用 CLIP ViT-B/16，关系用 all-MiniLM-L6-v2，属于评测编码器，与部件检测时 RN50 特征无关。脚本计算轴对齐包围盒体积 IoU/IoP，并使用严格大于门限的判断；这与第 4 节邻域覆盖 H/G 完全不同。



设真值节点为 $g$，预测为 $p$，$r_g(p)$ 为真值标签在固定词表中按文本相似度检索的名次，主节点协议可写为：



\[
R_{\mathrm{node}}@K=\frac{1}{|V_{\mathrm{GT}}|}
\sum_{g\in V_{\mathrm{GT}}}
\mathbf1\!\left[\exists p:\ \operatorname{IoU}_{3D}(g,p)>0,\ r_g(p)\le K\right].
\]

三元组还要求两个端点和关系标签同时检索命中。因此 R@5 不是“85\% 的实例身份一定正确”，也不是整张图只有五个候选。原脚本逐个真值寻找任意预测，\textbf{没有全局一对一匹配}；多余预测不会像 AP 那样直接被罚。



此外该脚本的 SPATIAL triplet 只评 \texttt{part of}，关系端点允许互换，不能据此评价全部有向空间边。层级功能图论文明确增加 Hungarian 一对一评测以避免一个预测命中多个真值；我们另列这种严格对应检查，保留 \texttt{part\_\allowbreak{}of} 子到父方向，不能把修改后的分数冠为未改动的 OP3DSG 指标。固定词表、标签别名、坐标系、几何尺度、匹配规则及评测版本后再比较方法，禁止为每个方法临时修改词表或阈值。



\subsection*{10.3 本地真值与当前可评范围}
\addcontentsline{toc}{subsection}{10.3 本地真值与当前可评范围}

已确认 \texttt{/\allowbreak{}home/\allowbreak{}lewisliu/\allowbreak{}datasets/\allowbreak{}scannet/\allowbreak{}scans/\allowbreak{}scene0000\_\allowbreak{}00} 存在 \texttt{aggregation.\allowbreak{}json}、\texttt{segs.\allowbreak{}json}、标注网格和 2D 实例 ZIP，可以准备官方实例评测。预测需落到 \texttt{vh\_\allowbreak{}clean\_\allowbreak{}2.\allowbreak{}ply} 的相同顶点顺序，并提供独立的实例置信排序；当前代表观测分数并非校准的轨迹置信度，不能未经记录便当作新的精度结论。



现有 30 帧实验属于部分场景观测。全场景官方 AP 会计入未见物体造成的漏检；只评可见区域是另一个协议，应另列并固定可见性规则。一个场景的试评也不能称为 ScanNet 官方测试集成绩。



用户指定 Hypersim 输入目录目前未发现语义/实例 HDF5 真值；已导出的深度和位姿不是关系标注。ZED \texttt{scene0802\_\allowbreak{}00} 也没有官方实例、部件与关系真值。二者需补充原始标注或独立人工标注，特别是功放碎片、相邻同类物体、部件父归属。ScanNet 自带实例标签不能凭空提供 \texttt{part\_\allowbreak{}of} 或功能关系；用待评算法的规则生成“真值”会构成循环验证。



\subsection*{10.4 实验与贡献判断}
\addcontentsline{toc}{subsection}{10.4 实验与贡献判断}

固定同一帧清单、矫正位姿、词表和真值，按场景隔离开发与评测；调参只用开发场景。先建立 DINO/SAM+legacy 与 Florence/SAM+legacy 的物体对照，再在同一 Florence 前端下分别加入部件层、改进关联和软父归属。部件层前后物体字段相同，应只主张部件/关系表达的变化，不能把其作为物体融合收益。颜色、SAM、去噪再单独消融；新增层级时，同时报告相同粒度子图与完整新增层级，避免分母改变。



报告分场景指标、跨场景汇总与耗时；AP 保留官方聚合，图检索记录原脚本计数及相应分母，另列场景均值。重复/误合并案例作为误差分析。没有足够场景时标注为试验结果；未预测或无真值的功能指标标为不可评，并说明原因。当前 scene0802 两路线同时改变前端和部件层，不能据其节点数归因于关联算法。上述协议已确定，但本次只完成文献、源码与数据核查，未执行评测或生成新实验产物。


\end{document}
